\documentclass[12pt]{book}
\usepackage{precalc1}
\setcounter{chapter}{8}
\begin{document}

\chapter{Systems of Linear Equations}

Until now, most of our attention has gone to one rule at a time: one function, one
formula, one graph. A system of equations changes the question. Instead of asking
what one equation says, we ask where several equations are true at the same time.
That is why systems have both an algebraic life and a geometric life. Algebraically,
we are solving several conditions together. Geometrically, we are looking for an
intersection.

A single linear equation in two variables describes a line; a pair of them asks
where two lines meet. A linear equation in three variables describes a plane; a
system of them asks where several planes meet. The solution set of a linear system
is the intersection of these flat objects, and that intersection takes one of exactly
three forms: a single point, nothing at all, or an entire shared line or plane. The
geometry tells us what kind of answer is possible, and the algebra finds the answer
when one exists.

\section{Two Equations in Two Variables}

Each equation $ax+by=c$ graphs as a line in the plane. A solution of the system is a
point lying on both lines at once, so the solution set is the intersection of the two
lines. This is the central translation: to solve the system is to find the point, or
points, shared by the graphs.

Two lines in a plane have very little freedom. They can cross once, never meet, or be
the same line. Those three pictures become the three possible outcomes for the
system.

\begin{keyidea}
\textbf{Two lines in the plane.}
\begin{itemize}\setlength{\itemsep}{2pt}
  \item \emph{Crossing} at one point: a \emph{unique solution}. The lines have
        different slopes.
  \item \emph{Parallel} and distinct: \emph{no solution}. Same slope, different
        intercepts; the lines never meet.
  \item \emph{Coincident} (the same line): \emph{infinitely many solutions}. Every
        point of the shared line solves both equations.
\end{itemize}
\end{keyidea}

\begin{center}
\begin{tikzpicture}[scale=0.62]
  % unique
  \begin{scope}
    \draw[->] (-2.2,0)--(2.6,0); \draw[->] (0,-2.2)--(0,2.6);
    \draw[pcblue,thick] (-2,-1.2)--(2.4,2.0);
    \draw[pcred,thick] (-2,1.6)--(2.4,-1.7);
    \fill (0.35,0.55) circle (2.2pt);
    \node at (0,-3.2) {\scriptsize unique solution};
  \end{scope}
  % none
  \begin{scope}[xshift=6cm]
    \draw[->] (-2.2,0)--(2.6,0); \draw[->] (0,-2.2)--(0,2.6);
    \draw[pcblue,thick] (-2.2,-1.0)--(2.4,1.3);
    \draw[pcred,thick] (-2.2,0.2)--(2.4,2.5);
    \node at (0,-3.2) {\scriptsize no solution};
  \end{scope}
  % infinite
  \begin{scope}[xshift=12cm]
    \draw[->] (-2.2,0)--(2.6,0); \draw[->] (0,-2.2)--(0,2.6);
    \draw[pcblue,thick] (-2.2,-1.1)--(2.4,1.2);
    \draw[pcred,thick,dashed] (-2.0,-1.0)--(2.2,1.1);
    \node at (0,-3.2) {\scriptsize infinitely many};
  \end{scope}
\end{tikzpicture}
\end{center}

\begin{keyidea}
A linear system has exactly one of three outcomes: a unique solution, no solution, or
infinitely many. No linear system has, say, exactly two solutions; flat objects meet
in a flat object, and the only flat possibilities are a point, the empty set, or a
whole line or plane.
\end{keyidea}

This fact is useful before any calculation begins. When we solve a system, we are not
hunting for a mysterious fourth possibility. The algebra will eventually announce one
of the three outcomes above.

\section{Solving by Substitution}

Substitution is the most direct algebraic method for a small system. The idea is that
an equation gives permission to replace one expression by another equal expression.
If one equation says $y=7-2x$, then anywhere the other equation contains $y$, we may
write $7-2x$ instead. That replacement reduces two equations in two unknowns to one
equation in one unknown.

After the single equation is solved, the value is fed back into one of the original
equations to recover the other variable. The method is slow for large systems, but it
is excellent for seeing what a solution means.

\begin{example}[a unique solution by substitution]
Solve
\[
\begin{aligned}
2x+y&=7,\\
3x-2y&=0.
\end{aligned}
\]
This system should have a single solution if the two lines cross. The substitution
will find the coordinates of that crossing point.

Isolate $y$ in the first equation: $y=7-2x$. Substitute into the second:
$3x-2(7-2x)=0$, so $3x-14+4x=0$, giving $7x=14$ and $x=2$. Back-substitute:
$y=7-2(2)=3$. The unique solution is $(2,3)$, the crossing point of the two lines.
Check: $2(2)+3=7$ and $3(2)-2(3)=0$.
\end{example}

\begin{example}[detecting no solution]
Solve
\[
\begin{aligned}
x-y&=1,\\
x-y&=4.
\end{aligned}
\]
Here the two left sides are identical, but the right sides disagree. The system is
asking the same expression to equal two different numbers at once.

Substituting $x=y+1$ from the first equation into the second gives $(y+1)-y=4$, that
is $1=4$, a false statement. A contradiction signals that the lines are parallel and
distinct: there is no solution.
\end{example}

\begin{example}[detecting infinitely many]
Solve
\[
\begin{aligned}
x+2y&=3,\\
2x+4y&=6.
\end{aligned}
\]
The second equation is twice the first, so it carries no new information. The two
lines are really the same line written in different forms.

Substituting $x=3-2y$ into the second equation gives $2(3-2y)+4y=6$, that is $6=6$,
always true. Every point on $x+2y=3$ solves the system; the solution set is that
entire line. Writing $y=t$ as a free parameter, the solutions are $(3-2t,\,t)$ for
all $t$.
\end{example}

\begin{keyidea}
Substitution reveals the outcome through the equation it produces. A determined value
gives a unique solution; a false statement such as $1=4$ signals no solution; an
identity such as $6=6$ signals infinitely many, described by a free parameter.
\end{keyidea}

The false statement and the identity are not failures of the method. They are the
method reporting the geometry: parallel lines give a contradiction, and coincident
lines give an equation that is automatically true.

\section{Three Equations in Three Variables}

A linear equation $ax+by+cz=d$ in three variables graphs as a plane in space. A
system of three such equations asks for the points common to three planes, and the
solution set is their intersection. The same three outcomes appear, now staged by
planes instead of lines.

The pictures are harder to draw accurately, but the logic has not changed. A unique
solution is one common point. No solution means the planes have no point shared by
all of them. Infinitely many solutions means the planes share a whole line or a whole
plane.

\begin{keyidea}
\textbf{Three planes in space.}
\begin{itemize}\setlength{\itemsep}{2pt}
  \item \emph{Unique solution:} the three planes meet at a single point, like three
        walls meeting at a corner.
  \item \emph{No solution:} the planes have no common point, for instance two
        parallel planes, or three forming an open triangular prism that never closes.
  \item \emph{Infinitely many:} the planes share a whole line, or all coincide in a
        single plane.
\end{itemize}
\end{keyidea}

\begin{center}
\begin{tikzpicture}[scale=0.95]
  % three planes meeting at a corner (unique)
  \begin{scope}
    % floor
    \fill[pcblue!20] (-1.2,-0.6)--(1.6,-0.6)--(1.0,0.5)--(-1.8,0.5)--cycle;
    % back wall
    \fill[pcred!20] (-1.2,-0.6)--(-1.8,0.5)--(-1.8,2.1)--(-1.2,1.0)--cycle;
    % side wall
    \fill[pcgreen!25] (-1.2,-0.6)--(1.6,-0.6)--(1.6,1.0)--(-1.2,1.0)--cycle;
    \fill (-1.2,-0.6) circle (1.8pt);
    \node at (0,-1.4) {\scriptsize unique: a corner};
  \end{scope}
  % parallel planes (none)
  \begin{scope}[xshift=4.7cm]
    \fill[pcblue!25] (-1.6,0.2)--(1.2,0.2)--(0.6,1.2)--(-2.2,1.2)--cycle;
    \fill[pcred!25] (-1.6,-0.9)--(1.2,-0.9)--(0.6,0.1)--(-2.2,0.1)--cycle;
    \node at (-0.5,-1.7) {\scriptsize none: parallel};
  \end{scope}
  % common line (infinite)
  \begin{scope}[xshift=9.4cm]
    \fill[pcblue!20] (-1.4,-0.8)--(0.2,-1.1)--(0.9,1.0)--(-0.7,1.3)--cycle;
    \fill[pcgreen!22] (0.2,-1.1)--(1.6,-0.6)--(0.9,1.4)--(-0.5,0.9)--cycle;
    \draw[pcpurple,very thick] (-0.25,-0.95)--(0.5,1.2);
    \node at (0,-1.8) {\scriptsize infinite: shared line};
  \end{scope}
\end{tikzpicture}
\end{center}

\section{Back-Substitution and Triangular Form}

A three-variable system is easiest to solve when it is in \emph{triangular form}: the
first equation involves all three variables, the second only two, the third only
one. The shape of the system has become a staircase. The last equation gives one
variable immediately; substituting upward (back-substitution) recovers the others one
at a time.

This is the same idea as working backward through a set of instructions. Once the
last variable is known, the equation just above it becomes simple. Once two variables
are known, the top equation becomes simple.

\begin{keyidea}
A system is in \emph{triangular form} when each successive equation drops a variable,
leaving one equation in a single unknown. Solve that equation, then
\emph{back-substitute} the value upward to peel off the remaining variables in turn.
\end{keyidea}

\begin{example}[back-substitution in a triangular system]
Solve
\[
\begin{aligned}
x+2y-z&=3,\\
y+2z&=4,\\
3z&=9.
\end{aligned}
\]
The third equation gives $z=3$. Substitute into the second: $y+6=4$, so $y=-2$.
Substitute both into the first: $x+2(-2)-3=3$, so $x-7=3$ and $x=10$. The unique
solution is $(10,-2,3)$. Each step uses only values already found, moving upward.
\end{example}

The work of solving a general system is to reach triangular form by eliminating
variables. Adding a multiple of one equation to another removes a chosen variable
without changing the solution set, because the new equation holds exactly when the
old ones do. We are not changing the problem; we are rewriting it into a form whose
solution can be read from the bottom upward.

\begin{example}[reducing to triangular form, then back-substituting]
Solve
\[
\begin{aligned}
x+y+z&=6,\\
2x-y+z&=3,\\
x+2y-z&=2.
\end{aligned}
\]
Subtract twice the first equation from the second:
$(2x-y+z)-2(x+y+z)=3-12$, giving $-3y-z=-9$. Subtract the first equation from the
third: $y-2z=-4$. The system is now
\[
\begin{aligned}
x+y+z&=6,\\
-3y-z&=-9,\\
y-2z&=-4.
\end{aligned}
\]
From the last two equations, eliminate $y$: multiply the third by $3$ to get
$3y-6z=-12$ and add to the second, giving $-7z=-21$, so $z=3$. Back-substitute into
$y-2z=-4$: $y=2$. Then the first equation gives $x=6-2-3=1$. The solution is
$(1,2,3)$. Check in all three originals: $1+2+3=6$, $2-2+3=3$, $1+4-3=2$.
\end{example}

\begin{pitfall}
Elimination changes the equations but never the solution set, provided each step is
reversible: adding a multiple of one equation to another, or scaling an equation by a
nonzero constant. Scaling by zero is not allowed, since it discards an equation and
can manufacture solutions the original system never had.
\end{pitfall}

\section{An Application}

Applications of systems usually begin as plain-language constraints. Each sentence
becomes an equation, and each unknown represents a quantity we need to determine. The
important habit is to name the variables first, then translate each condition without
trying to solve too early.

\begin{example}[mixing for a target]
A chemist combines three solutions of $10\%$, $20\%$, and $50\%$ concentration to
make $1$ liter of $22\%$ solution, using twice as much of the $10\%$ as the $50\%$.
Let $x,y,z$ be the liters of the $10\%$, $20\%$, $50\%$ solutions. The conditions are
\[
\begin{aligned}
x+y+z&=1 &&\text{(total volume)},\\
0.10x+0.20y+0.50z&=0.22 &&\text{(total active ingredient)},\\
x-2z&=0 &&\text{(twice as much }10\%\text{ as }50\%).
\end{aligned}
\]
From the third equation, $x=2z$. Substituting into the first two equations and solving
the resulting two-variable system gives $z=0.2$, $x=0.4$, and $y=0.4$ liters. Three
plain-language conditions become three linear equations, and their common solution is
the recipe.
\end{example}

The point is not that every word problem is secretly about chemistry. The point is
that a system is a way of making several requirements hold at once: total amount,
total concentration, and an additional relationship among the ingredients.

\section{Exercises}

\begin{exercises}
\item Describe the geometry and the number of solutions for each:
(a) two lines with different slopes; (b) two distinct lines with equal slope;
(c) one line written two equivalent ways.

\item Solve by substitution: $\begin{aligned}x+y&=5\\ 2x-y&=1\end{aligned}$ and check
both equations.

\item Solve by substitution and identify the outcome:
$\begin{aligned}3x-y&=2\\ 6x-2y&=10\end{aligned}$.

\item Solve and identify the outcome:
$\begin{aligned}x-2y&=4\\ -2x+4y&=-8\end{aligned}$, writing the solution set with a
free parameter if appropriate.

\item Explain why a linear system can never have exactly two solutions, in terms of
the geometry of lines and planes.

\item Describe two distinct geometric arrangements of three planes that each yield no
solution.

\item Solve the triangular system by back-substitution:
$\begin{aligned}x-y+2z&=5\\ y-z&=1\\ 4z&=8\end{aligned}$.

\item Reduce to triangular form and solve:
$\begin{aligned}x+y+z&=4\\ 2x+y-z&=3\\ x-2y+3z&=11\end{aligned}$, and check your
solution.

\item Solve and classify:
$\begin{aligned}x+y+z&=2\\ 2x+2y+2z&=4\\ x-y+z&=0\end{aligned}$, describing the
solution set geometrically.

\item A theater sells adult tickets at \$12 and child tickets at \$8. One evening
$200$ tickets bring in \$2000. Set up and solve a system for the numbers of each
kind.

\item A shop blends coffee at \$8, \$10, and \$14 per pound into $10$ pounds selling
at \$11 per pound, using equal amounts of the two cheaper beans. Write and solve the
system for the three weights.

\item Explain why adding a multiple of one equation to another leaves the solution
set unchanged, and why scaling an equation by zero is forbidden.
\end{exercises}

\end{document}
